% SGA 3, Expose VI_A, Loop 1 English reconstruction.
% Scope: paragraph 3.2.2 in full; paragraph 3.2.3 is excluded.
% Authority: Polo--Gille Exp6A-13oct24.pdf, component-local page 23
% (combined-reader page 326; printed page 315).
% OCR and the Reinhold English Markdown are comparison witnesses only.

\medskip
\phantomsection\label{par:sga3-VIA-3.2.2}%
\noindent\textbf{3.2.2.---}
对每个为局部环且在 \(A\) 上有限自由的 \(A\)-代数 \(A'\),
用 \(U(A')\) 表示 \(G\otimes_AA'\) 中具有饱和开邻域 \(W\) 的那些点组成的集合,
使 \(G_*\otimes_AA'\) 在 \(W\) 上诱导的群胚有拟截面.
显然 \(U(A')\) 对 \(G(\operatorname{Spec}A')\) 在 \(G\otimes_AA'\) 上的作用饱和.
我们将证明, 当 \(A'\) 充分大时,
\[
  U(A')=G\otimes_AA'.
\]

由\SGARefLink{par:sga3-VIB-8.1}{定理~V.8.1}, \(U(A)\) 非空,
从而包含闭点 \(y\). 对下述维数作归纳:
\[
  \dim\bigl(G-U(A)\bigr).
\]
设 \(g_1,\ldots,g_n\) 为分别属于 \(G-U(A)\) 各不可约分支的闭点.
由\SGARefLink{prop:sga3-VIA-1.2}{1.2},
存在为局部环且在 \(A\) 上有限自由的 \(A\)-代数 \(A'\), 使得
\[
  g'_1,\ldots,g'_p
  \qquad
  \bigl(\text{相应地, }x=x_1,\ldots,x_r\bigr)
\]
为 \(G'=G\otimes_AA'\) 中位于 \(g_1,\ldots,g_n\)
上方 (相应地, 位于 \(y\) 上方) 的点, 且在 \(A'\) 上严格有理.
于是 \(U(A')\) 包含
\[
  \bigl(U(A)\otimes_AA'\bigr)\cdot x^{-1}g'_i
  \qquad\text{对每个 }i.
\]
因此 \(U(A')\) 包含 \(g'_1,\ldots,g'_p\), 且
\[
  \dim\bigl(G'-U(A')\bigr)
  <
  \dim\bigl(G-U(A)\bigr).
\]

因此由归纳假设得到一个为局部环且在 \(A'\) 上有限自由的代数 \(A''\), 使得
\[
  U(A'')
  =
  G'\otimes_{A'}A''
  =
  G\otimes_AA''.
\]
% Independent prose review: source and translation compared paragraph by paragraph through this component; statements, hypotheses, proofs, notes and diagram captions covered; bibliography reserved for its assigned reviewer; source math and reference keys unchanged.
